\documentclass[11pt]{article}
\usepackage{mathblatt}
\usepackage{multicol}


\begin{document}
\pfheftstil
\pfheftkopf{Kurvenuntersuchung \textperiodcentered{} Lösungen}{}
\begin{pfloesung}{}
\lz{1.}{Punkte eintragen: $(-2 | -2)$, $(-1 | 2)$, $(0 | 0)$, $(1 | -2)$, $(2 | 2)$}{$f(-2) = -2$, $f(-1) = 2$, $f(0) = 0$, $f(1) = -2$, $f(2) = 2$}{}
\end{pfloesung}
\begin{pfloesung}{}
\lz{2.}{$x = 3$}{$A'(x) = 27 - 3x^2$; $27 - 3x^2 = 0 \Rightarrow x^2 = 9$}{}
\end{pfloesung}
\begin{pfloesung}{}
\lz{3.}{Faktoren null setzen}{}{}
\end{pfloesung}
\begin{pfloesung}{}
\lz{4.}{Stelle gegeben, Wert gesucht}{}{}
\end{pfloesung}
\begin{pfloesung}{}
\lz{5.}{f(1) = 2e $-$ 2e = 0}{}{}
\end{pfloesung}
\begin{pfloesung}{}
\lz{6.}{80 °C}{$k(0) = 60 + 20$}{}
\end{pfloesung}
\begin{pfloesung}{}
\lz{7.}{x$_1$ = 0, x$_2$ = 2, x$_3$ = 6}{$f(x) = x(x - 2)(x - 6)$}{}
\end{pfloesung}
\begin{pfloesung}{}
\lz{8.}{e$^{x}$ > 0 für alle x $\Rightarrow$ 3e$^{x}$ + 1 > 1 > 0 $\Rightarrow$ keine Nullstelle; Wertemenge ]1; $\infty$[}{}{}
\end{pfloesung}
\begin{pfloesung}{{\small $f'(x) = 3x^{2} - 24x + 45$}\par\noindent{\small $f''(x) = 6x - 24$}\par\noindent{\small $g'(x) = 1{,}2(x - 5)^{2}$}}
\lz{9.}{S$_y$(0 | $-$50); genau zwei Nullstellen x$_1$ = 2 und x$_2$ = 5 (doppelt)}{$f(x) = (x - 5)^{2} \cdot  (x - 2)$}{}
\end{pfloesung}
\begin{pfloesung}{{\small $f''(x) = 6x^{2} - 8$}}
\lz{10.}{N$_1$($-$1 | 0) wegen Achsensymmetrie; weitere Schnittpunkte N$_2$($-\sqrt{7}$ | 0), N$_3$($\sqrt{7}$ | 0)}{mit a); u = x²: 0,5u² $-$ 4u + 3,5 = 0 $\Rightarrow$ u = 1 oder u = 7}{}
\end{pfloesung}
\begin{pfloesung}{}
\lz{11.}{Nullstelle x = 2; S$_y$(0 | $-$2e³) $\approx$ (0 | $-$40,2)}{$f(0) = -2e^{3} \approx  -40{,}17$}{}
\end{pfloesung}
\begin{pfloesung}{{\small $f''(x) = (0{,}5x^{2} + 2x - 1) \cdot  e^x$}}
\lz{12.}{S$_y$(0 | $-$2), N$_1$($-$2 | 0), N$_2$(2 | 0); R(0 | 2)}{$f(0) = -2$; $x^{2} - 4 = 0$}{}
\end{pfloesung}
\begin{pfloesung}{{\small $h'(t) = \tfrac{27}{40} t^{2} - 27t + \tfrac{405}{2}$}\par\noindent{\small $h''(t) = \tfrac{27}{20} t - 27$}}
\lz{13.}{N($-$2 | 0), S$_y$(0 | 2); f'(x) = 0 $\Rightarrow$ x = $-$1, Vorzeichenwechsel von + nach $-$ $\Rightarrow$ H($-$1 | e)}{$f(0) = 2$; $f(-1) = e \approx  2{,}72$}{}
\end{pfloesung}
\begin{pfloesung}{}
\lz{14.}{nachweisen: $x = 2$ einsetzen}{}{}
\end{pfloesung}
\begin{pfloesung}{}
\lz{15.}{gefragt ist nur, wo die Tangente waagerecht ist}{}{}
\end{pfloesung}
\begin{pfloesung}{}
\lz{16.}{x$_1$ = 0 und x$_2$ = 2}{$f'(x) = 3x^{2} - 6x$}{}
\end{pfloesung}
\begin{pfloesung}{{\small $f'(x) = (1 - x) \cdot  e^x$}}
\lz{17.}{H(1 | e)}{$f'(x) = 0 \Rightarrow  x = 1$}{}
\end{pfloesung}
\begin{pfloesung}{{\small $f'(x) = -\tfrac{1}{2} x^{2} + x$}\par\noindent{\small $f''(x) = -x + 1$}}
\lz{18.}{T(0 | 0), H(2 | $\tfrac{2}{3}$)}{}{}
\end{pfloesung}
\begin{pfloesung}{{\small $f''(x) = \tfrac{12}{25} x$}}
\lz{19.}{H($-$2,5 | 2,5), T(2,5 | $-$2,5); Gerade durch beide: y = $-$x $\Rightarrow$ Winkelhalbierende des II. und IV. Quadranten}{$f'(x) = \tfrac{6}{25} x^{2} - \tfrac{3}{2} = 0 \Rightarrow  x = \pm 2{,}5$}{}
\end{pfloesung}
\begin{pfloesung}{}
\lz{20.}{$y = -x$}{$f'(x) = \tfrac{3}{8}x^2 - 1{,}5 = 0 \Rightarrow x^2 = 4$; $f''(x) = \tfrac{3}{4}x$: $H(-2 | 2)$, $T(2 | -2)$}{}
\end{pfloesung}
\begin{pfloesung}{}
\lz{21.}{h'(10) = 0, h''(10) = $-$13,5 < 0 $\Rightarrow$ Maximum; h(10) = 900 m}{Kopf bei Nr. 13}{}
\end{pfloesung}
\begin{pfloesung}{}
\lz{22.}{f'($-$1 + $\sqrt{5}$) = 0 $\Rightarrow$ waagerechte Tangente; f($-$1 + $\sqrt{5}$) = (1 $-$ $\sqrt{5}$) · e$^{√5 - 1}$ $\approx$ $-$4,25}{Kopf bei Nr. 12; $(-1 + \sqrt{5})^{2} = 6 - 2\sqrt{5}$}{}
\end{pfloesung}
\begin{pfloesung}{}
\lz{23.}{H(3 | 4), T(5 | 0)}{Kopf bei Nr. 9}{}
\end{pfloesung}
\begin{pfloesung}{}
\lz{24.}{H(0 | 3,5); T$_1$($-$2 | $-$4,5), T$_2$(2 | $-$4,5)}{Kopf bei Nr. 10; $f'(x) = 2x^{3} - 8x = 2x(x - 2)(x + 2)$}{}
\end{pfloesung}
\begin{pfloesung}{}
\lz{25.}{H(4 | 2e) $\approx$ (4 | 5,44)}{$f'(x) = 0 \Rightarrow  x = 4$; Vorzeichenwechsel von + nach $-$}{}
\end{pfloesung}
\begin{pfloesung}{}
\lz{26.}{der Wert $f'''(1) = 6$ ist gegeben}{}{}
\end{pfloesung}
\begin{pfloesung}{}
\lz{27.}{gefragt ist die größte Wachstumsrate}{}{}
\end{pfloesung}
\begin{pfloesung}{}
\lz{28.}{W(4 | 2)}{Kopf bei Nr. 9}{}
\end{pfloesung}
\begin{pfloesung}{}
\lz{29.}{Wendepunkt (t = 20): Zeitpunkt der größten Geschwindigkeit beim Aufstieg}{Kopf bei Nr. 13; $h''(t) = 0 \Rightarrow  t = 20$}{}
\end{pfloesung}
\begin{pfloesung}{}
\lz{30.}{a = 0 (Tiefpunkt bei 0,2); b $\approx$ 0,5 (jede Wendestelle, auch 1,2 oder 1,8)}{}{}
\end{pfloesung}
\begin{pfloesung}{}
\lz{31.}{P $\approx$ (1,8 | 0,22); Wendepunkt (Krümmungswechsel) im fallenden Abschnitt zwischen Hochpunkt ($\approx$ 1,6) und Tiefpunkt ($\approx$ 2) $\Rightarrow$ f''(x\_P) = 0 und f'(x\_P) < 0}{mit a)}{}
\end{pfloesung}
\begin{pfloesung}{}
\lz{32.}{f''(x) = (x/4 $-$ $\tfrac{3}{2}$) · e$^{-x/2 + 3}$ mit Produktregel; f''(6) = 0 mit Vorzeichenwechsel, f(6) = 4 $\Rightarrow$ W(6 | 4)}{$f'''(6) = \tfrac{1}{4}$}{}
\end{pfloesung}
\begin{pfloesung}{}
\lz{33.}{W1($-$2 | $\tfrac{2}{3}$), W2(0 | 2)}{$f''(x) = x^{2} + 2x = x(x + 2)$; $f'''(x) = 2x + 2: f'''(0) = 2, f'''(-2) = -2$}{}
\end{pfloesung}
\begin{pfloesung}{}
\lz{34.}{W1($-$4,45 | 0,09), W2(0,45 | $-$2,98)}{Kopf bei Nr. 12; $x = -2 \pm  \sqrt{6}$}{}
\end{pfloesung}
\begin{pfloesung}{}
\lz{35.}{die Tangente in $P$ geht nach rechts oben}{}{}
\end{pfloesung}
\begin{pfloesung}{}
\lz{36.}{$A$: plus, $B$: null, $C$: minus}{}{}
\end{pfloesung}
\begin{pfloesung}{}
\lz{37.}{f' wechselt an seiner einzigen Nullstelle $-$3 das Vorzeichen von $-$ nach + $\Rightarrow$ f fällt für x $\le$ $-$3 und steigt für x $\ge$ $-$3}{}{}
\end{pfloesung}
\begin{pfloesung}{}
\lz{38.}{Graph II ist f'; Nullstelle bei x = $-$1 (Hochstelle von f), negativ für x > $-$1, weil f dort fällt}{Kopf bei Nr. 13; mit a)}{}
\end{pfloesung}
\begin{pfloesung}{}
\lz{39.}{Q(2 | $-\tfrac{4}{3}$): tiefster Punkt des unteren Schienenverlaufs rechts, am Übergang in den rechten Halbkreis}{$g'(x) = x^{2} + 2x = 0 \Rightarrow  x = -2 (x \neq  0) \Rightarrow  x_Q = 2$; $y_Q = -g(-2) = -\tfrac{4}{3}$}{}
\end{pfloesung}
\begin{pfloesung}{{\small $f'(x) = 0{,}5 \cdot  e^{0{,}5x}$}}
\lz{40.}{f'(x) = 0,5 · e$^{0{,}5x}$ > 0 für alle x $\Rightarrow$ streng monoton steigend}{}{}
\end{pfloesung}
\begin{pfloesung}{}
\lz{41.}{genau eine Nullstelle von g' (x = 5) $\Rightarrow$ genau eine waagerechte Tangente; g' $\ge$ 0 ohne Vorzeichenwechsel $\Rightarrow$ kein Extrempunkt (Sattelpunkt)}{Kopf bei Nr. 9}{}
\end{pfloesung}
\begin{pfloesung}{}
\lz{42.}{streng monoton fallend auf [0; $-$1 + $\sqrt{5}$], streng monoton steigend auf [$-$1 + $\sqrt{5}$; $\infty$)}{Kopf bei Nr. 12; mit c); Tiefstelle $-$1 + $\sqrt{5}$ $\approx$ 1,24}{}
\end{pfloesung}
\begin{pfloesung}{}
\lz{43.}{I falsch}{Kopf bei Nr. 12; mit c); Tiefpunkt hat den kleinsten Wert seiner Umgebung $\Rightarrow$ f($-$0,9 + $\sqrt{5}$) > f($-$1 + $\sqrt{5}$). II wahr; $f(-0{,}9 + \sqrt{5}) \approx  -4{,}21$; $f'(-0{,}9 + \sqrt{5}) \approx  0{,}87$}{}
\end{pfloesung}
\begin{pfloesung}{}
\lz{44.}{Grad ungerade, Koeffizient negativ}{}{}
\end{pfloesung}
\begin{pfloesung}{}
\lz{45.}{der Faktor $e^{-3x}$}{}{}
\end{pfloesung}
\begin{pfloesung}{}
\lz{46.}{Grenzwert 0; Graph nähert sich der x-Achse von oben}{Kopf bei Nr. 13}{}
\end{pfloesung}
\begin{pfloesung}{}
\lz{47.}{x $\to$ +$\infty$: f(x) $\to$ $-\infty$; x $\to$ $-\infty$: f(x) $\to$ +$\infty$}{Kopf bei Nr. 18; Leitterm $-\tfrac{1}{6}$ x³}{}
\end{pfloesung}
\begin{pfloesung}{}
\lz{48.}{x $\to$ $-\infty$: f'(x) $\to$ $-\infty$; x $\to$ +$\infty$: f'(x) $\to$ +$\infty$}{Leitterm $\tfrac{1}{3}$ x³}{}
\end{pfloesung}
\begin{pfloesung}{}
\lz{49.}{Nullstelle 2; x $\to$ $-\infty$: f(x) $\to$ 0; x $\to$ +$\infty$: f(x) $\to$ $-\infty$}{Kopf bei Nr. 17}{}
\end{pfloesung}
\begin{pfloesung}{}
\lz{50.}{x $\to$ +$\infty$: f(x) $\to$ +$\infty$; x $\to$ $-\infty$: f(x) $\to$ 0}{Kopf bei Nr. 12}{}
\end{pfloesung}
\begin{pfloesung}{}
\lz{51.}{achsensymmetrisch}{}{}
\end{pfloesung}
\begin{pfloesung}{}
\lz{52.}{nur ungerade Exponenten $\Rightarrow$ punktsymmetrisch zum Ursprung}{}{}
\end{pfloesung}
\begin{pfloesung}{}
\lz{53.}{nur gerade Exponenten (f($-$x) = f(x)) $\Rightarrow$ achsensymmetrisch zur y-Achse}{Kopf bei Nr. 10}{}
\end{pfloesung}
\begin{pfloesung}{}
\lz{54.}{f($-$x) = $-\tfrac{2}{25}$ x³ + $\tfrac{3}{2}$ x = $-$f(x) $\Rightarrow$ punktsymmetrisch; Nullstellen 0, $\pm$5$\sqrt{3}$/2 $\approx$ $\pm$4,33}{Kopf bei Nr. 19; $f(x) = x \cdot  (\tfrac{2}{25} x^{2} - \tfrac{3}{2})$; $x^{2} = \tfrac{75}{4}$}{}
\end{pfloesung}
\begin{pfloesung}{}
\lz{55.}{innen: x-Richtung}{}{}
\end{pfloesung}
\begin{pfloesung}{}
\lz{56.}{$4$, das Plus heißt nach links}{$2x + 8 = 2 \cdot (x + 4)$}{}
\end{pfloesung}
\begin{pfloesung}{}
\lz{57.}{falsch}{mit a); Verschiebung in y-Richtung erhält Steigungen, aber g'(0) = 2 $\neq$ h'(0) = 1; $h'(0) = 1$}{}
\end{pfloesung}
\begin{pfloesung}{}
\lz{58.}{Spiegelung an der y-Achse und Streckung in x-Richtung mit Faktor 400}{}{}
\end{pfloesung}
\begin{pfloesung}{}
\lz{59.}{c = 4, d = 55; Graph von f um 55 nach oben verschoben, dann in y-Richtung mit Faktor 4 gestreckt}{Kopf bei Nr. 9; $f(x) + 55 = x^{3} - 12x^{2} + 45x + 5$}{}
\end{pfloesung}
\begin{pfloesung}{}
\lz{60.}{h(x) = f(x + 2) = e$^{0{,}5x + 1}$ $-$ e}{Kopf bei Nr. 40; Nullstelle von f: x = 2}{}
\end{pfloesung}
\begin{pfloesung}{}
\lz{61.}{z. B. g(x) = f($-$x) = 0,5 · (x² $-$ 4) · e$^{-x}$ (IV. Quadrant) und h(x) = $-$f(x) = $-$0,5 · (x² $-$ 4) · e$^{x}$ (II. Quadrant); k(x) = $-$f($-$x) (I. Quadrant)}{Kopf bei Nr. 12}{}
\end{pfloesung}
\begin{pfloesung}{}
\lz{62.}{g* $\perp$ g durch den Ursprung $\Rightarrow$ Steigung $-$1/($\tfrac{1}{2}$) = $-$2 $\Rightarrow$ g*(x) = $-$2x; A*($-$2,5 | 5)}{Kopf bei Nr. 19; Drehung um 90°: (x | y) $\to$ ($-$y | x)}{}
\end{pfloesung}
\begin{pfloesung}{}
\lz{63.}{nein}{mit a); Graph von g entsteht aus Graph von f durch Streckung in x-Richtung (Faktor 3) und y-Richtung (Faktor 2), beide > 1 $\Rightarrow$ Strecke länger; $\mid AB\mid  = \sqrt{\pi ^{2} + 4} \approx  3{,}72$; bei g: $\sqrt{9 \pi ² + 16}$ $\approx$ 10,24}{}
\end{pfloesung}
\begin{pfloesung}{}
\lz{64.}{z = $\tfrac{25}{3}$ $\approx$ 8,33}{mit d); Steigungsdreieck 3 : 4 : 5; Abstand 5 ist Kathete, z Hypotenuse: z/5 = $\tfrac{5}{3}$}{}
\end{pfloesung}
\begin{pfloesung}{}
\lz{65.}{g(x) = $-$(x $-$ 2) · e$^{-x/2 + 3}$; Durchmesser $\approx$ 7,4 mm}{$f(11) = 9e^{-2{,}5} \approx  0{,}739$; $2 \cdot  0{,}739 LE = 1{,}48 LE = 0{,}74 cm$}{}
\end{pfloesung}
\begin{pfloesung}{}
\lz{66.}{A = $\tfrac{1}{2}$ · 2$\pi$ · 2 = 2$\pi$}{Grundseite AC = 2$\pi$; Höhe 2}{}
\end{pfloesung}
\begin{pfloesung}{}
\lz{67.}{180 l}{Kopf bei Nr. 12; mit b); Grundfläche 60 cm $\times$ 60 cm; Höhe 50 cm}{}
\end{pfloesung}
\begin{pfloesung}{}
\lz{68.}{passt; Höhe 2 · $\tfrac{4}{3}$ = $\tfrac{8}{3}$ m $\approx$ 2,67 m $\le$ 2,7 m, Breite 2 · (2 + $\tfrac{4}{3}$) = $\tfrac{20}{3}$ m $\approx$ 6,67 m $\le$ 6,7 m}{Halbkreise um N($\pm$2 | 0) mit Radius $\tfrac{4}{3}$ reichen bis x = $\pm\tfrac{10}{3}$}{}
\end{pfloesung}
\begin{pfloesung}{}
\lz{69.}{$\approx$ 1041 m}{Kopf bei Nr. 13; mit j); $h(15) = 759{,}375$; Aufstieg 10 bis 15 min: 900 $-$ 759,4 = 140,6 m}{}
\end{pfloesung}
\begin{pfloesung}{}
\lz{70.}{Länge $\approx$ 18,4 m; Zeit $\approx$ 5,25 min $\approx$ 5 min 15 s}{vier Bögen 4 · 2,5 = 10 m; Vollkreis mit Radius $\tfrac{4}{3}$: 2$\pi$ · $\tfrac{4}{3}$ $\approx$ 8,38 m}{}
\end{pfloesung}
\begin{pfloesung}{}
\lz{71.}{4,5 cm $\times$ 5,44 cm $\times$ 5,44 cm; Masse $\approx$ 114 g}{mit b); Länge 9 LE = 4,5 cm; Breite und Höhe 2 · 2e $\approx$ 10,87 LE $\approx$ 5,44 cm}{}
\end{pfloesung}
\begin{pfloesung}{}
\lz{72.}{Substitution}{}{}
\end{pfloesung}
\begin{pfloesung}{}
\lz{73.}{$x > 3$}{}{}
\end{pfloesung}
\begin{pfloesung}{}
\lz{74.}{$x_1 = 2$, $x_2 = -2$}{$f'(x) = -2x + 6$; $-x^2 + 6x - 4 = -2x^2 + 6x$, $x^2 = 4$}{}
\end{pfloesung}
\begin{pfloesung}{}
\lz{75.}{f'(x) = $-$2x + 4; x$_1$ = $-$1 oder x$_2$ = 1}{$-x^{2} + 4x - 1 = -2x^{2} + 4x \Rightarrow  x^{2} = 1$}{}
\end{pfloesung}
\begin{pfloesung}{}
\lz{76.}{(x + 2) · e$^{-x}$ = $-$(x + 1) · e$^{-x}$ $\Rightarrow$ x + 2 = $-$x $-$ 1 $\Rightarrow$ x = $-$1,5}{Kopf bei Nr. 13}{}
\end{pfloesung}
\begin{pfloesung}{}
\lz{77.}{L = (0; 5) $\cup$ (5; $\infty$)}{Kopf bei Nr. 9; $f(x) - g(x) = 0{,}6x(x - 5)^{2}$}{}
\end{pfloesung}
\begin{pfloesung}{}
\lz{78.}{Lösungsweg: |p(x) $-$ f(x)| $\le$ $\tfrac{3}{8}$ ansetzen; für x $\ge$ 1 ist f $\ge$ p $\Rightarrow$ 0,5x$^{4}$ $-$ 0,5x² $\le$ $\tfrac{3}{8}$; Randgleichung mit u = x² lösen (u = 1,5) $\Rightarrow$ 1 $\le$ x $\le$ $\sqrt{1{,}5}$ $\approx$ 1,22}{Kopf bei Nr. 10; mit g); $0{,}5x^{4} - 0{,}5x^{2} = \tfrac{3}{8} \Rightarrow  x^{2} = 1{,}5$}{}
\end{pfloesung}
\begin{pfloesung}{}
\lz{79.}{x $\approx$ 1,9 (grafisch: Stelle, an der der Graph 3 Einheiten weiter rechts 1000 höher liegt); Deutung: Zeitpunkt, zu dem die Anzahl der Likes um 1000 kleiner ist als drei Stunden später}{Kopf bei Nr. 17}{}
\end{pfloesung}
\begin{pfloesung}{}
\lz{80.}{nach unten}{}{}
\end{pfloesung}
\begin{pfloesung}{}
\lz{81.}{$(0 | 2)$}{}{}
\end{pfloesung}
\begin{pfloesung}{}
\lz{82.}{Punkte eintragen: $(-2 | -6)$, $(-1 | 0)$, $(0 | 0)$, $(1 | 0)$, $(2 | 6)$}{$f(-2) = -6$, $f(-1) = 0$, $f(0) = 0$, $f(1) = 0$, $f(2) = 6$}{}
\end{pfloesung}
\begin{pfloesung}{}
\lz{83.}{Anzahl der Likes nimmt zunächst immer stärker zu, ab einem Zeitpunkt immer schwächer, nähert sich 2500}{Kopf bei Nr. 17}{}
\end{pfloesung}
\begin{pfloesung}{}
\lz{84.}{$4{,}5$}{Randwerte $f(-3) = f(3)$}{}
\end{pfloesung}
\begin{pfloesung}{}
\lz{85.}{$\approx$ 230 °C (228 bis 235 °C vertretbar)}{Kopf bei Nr. 9; Schwankung zwischen $\approx$ 222 °C und 245 °C}{}
\end{pfloesung}
\begin{pfloesung}{}
\lz{86.}{falsch}{Kopf bei Nr. 9; mit b), c); Wendetangente trifft den Graphen nur in W; Wendetangente y = $-$3x + 14; $f(7) = 20$}{}
\end{pfloesung}
\begin{pfloesung}{}
\lz{87.}{$f(3) = 1$ und $f'(3) = 0$}{}{}
\end{pfloesung}
\begin{pfloesung}{}
\lz{88.}{drei Unbekannte $a$, $b$, $c$ und drei Gleichungen}{}{}
\end{pfloesung}
\begin{pfloesung}{}
\lz{89.}{Steigung (1,2 + 0,4)/(2 $-$ 0,4) = 1, y-Achsenabschnitt $-$0,8 $\Rightarrow$ g(x) = x $-$ 0,8; beide Punkte erfüllen die Gleichung}{Kopf bei Nr. 18}{}
\end{pfloesung}
\begin{pfloesung}{}
\lz{90.}{a = 1/32, b = $\tfrac{1}{8}$; k(0) = 0, k'(0) = 0 (von selbst), k(2) = 1, k'(2) = 1,5}{Kopf bei Nr. 18; Steigung RS = 1,5; $16a + 4b = 1, 32a + 4b = 1{,}5$}{}
\end{pfloesung}
\begin{pfloesung}{}
\lz{91.}{k(x) = 2x³ $-$ 28x² + 130x + 20}{Kopf bei Nr. 9; $d = 20, c = 130$; $125a + 25b = -450, 75a + 10b = -130$}{}
\end{pfloesung}
\begin{pfloesung}{}
\lz{92.}{die Randwerte $A(0)$ und $A(4)$}{}{}
\end{pfloesung}
\begin{pfloesung}{}
\lz{93.}{$60$}{}{}
\end{pfloesung}
\begin{pfloesung}{}
\lz{94.}{$x = 2$}{$A'(x) = 12 - 3x^2$; $12 - 3x^2 = 0 \Rightarrow x^2 = 4$}{}
\end{pfloesung}
\begin{pfloesung}{}
\lz{95.}{d(x) = 0,5x² $-$ 0,5x$^{4}$ maximal bei x = $\pm$1/$\sqrt{2}$ mit d = $\tfrac{1}{8}$}{Kopf bei Nr. 10; $d'(x) = x - 2x^{3} = 0 \Rightarrow  x = 0, \pm 1/\sqrt{2}$; $d(0) = d(\pm 1) = 0$}{}
\end{pfloesung}
\begin{pfloesung}{}
\lz{96.}{P($\sqrt{2}$ | ($\sqrt{2}$ + 2) · e$^{-√2}$) $\approx$ (1,41 | 0,83)}{Kopf bei Nr. 13; $A(u) = u \cdot  f(u) = u(u + 2)e^{-u}$; $A'(u) = (2 - u^{2})e^{-u} = 0 \Rightarrow  u = \sqrt{2}$}{}
\end{pfloesung}
\end{document}